\documentclass[UTF8,fontset=none,11pt,a4paper]{ctexart}
\usepackage{ipara-handbook}
\handbooksetup{NET-08 应用层服务语义}{408 · 计算机网络 NET-08}{DNS · HTTP · FTP · Mail}
\providecommand{\tightlist}{\setlength{\itemsep}{0pt}\setlength{\parskip}{0pt}}
\providecommand{\passthrough}[1]{#1}
\begin{document}
\handbooktitle{应用层服务语义}{把字节组织成名字、资源与操作}{408 · NET-08}
\section{本册定位}\label{0-ux672cux518cux5b9aux4f4d}

本 Topic 回答：端到端传输只提供 byte stream 或
datagram，应用怎样定义对象、名字、操作、响应和状态，使不同程序能够共同完成
DNS、Web、文件与邮件服务？

本册拥有 application architecture、DNS、HTTP semantics、FTP
与电子邮件的服务分工，以及 application protocol 对 transport service
的选择逻辑。

DHCP 虽以 UDP application protocol
形式运行，但在本项目中由\href{../04_IP地址_子网与分组转发/README.md}{IP
配置与转发}拥有，因为其主要输出是主机进入 IP 世界所需的启动配置。

\section{独立阅读词典：应用协议在定义什么}

\begin{handbooktable}{L{2.3cm} L{4.3cm} Y}
\toprule
术语 & 本册中的可操作定义 & 不要混成什么 \\
\midrule
application protocol & 约定对象、名字、消息字段、动作顺序和响应含义 & 不只是一个端口号 \\
client/server 与 P2P & 前者由 client 请求、server 提供；后者参与者可同时请求和提供 & 都不等于 TCP/UDP \\
resource / representation & resource 是 URI 标识的抽象对象；representation 是某时刻某格式的可传输表示 & response body 不是 resource 全部身份 \\
URI / URL & URI 是资源标识语法；URL 是带定位信息的常见 URI 形式 & 不等于 DNS 返回的 IP \\
domain / zone & domain 是名字树名称；zone 是权威管理者实际负责的数据边界 & 不等于物理网络区域 \\
resolver / authoritative & resolver 代客户端查询并缓存；authoritative server 对自己的 zone 给权威答案；TTL 是缓存时间预算 & resolver 不是权威数据 Owner \\
request/response 与邮件/文件协议 & request 表达意图，response 返回状态/表示；SMTP 传邮件，POP3/IMAP 访问邮箱，FTP 分离控制与数据 & 都不是 TCP ACK \\
\bottomrule
\end{handbooktable}

本册所说的“语义”就是一个具体问题的答案：客户端发出这个消息后，服务器应把
哪个对象当作目标、执行什么动作、返回什么状态，以及下一步是否允许继续。这样读
HTTP、DNS、邮件和 FTP 时，不会把字段格式误当成服务含义。

\section{根本问题：transport 不知道``这段 byte
意味着什么''}\label{1-ux6839ux672cux95eeux9898transport-ux4e0dux77e5ux9053ux8fd9ux6bb5-byte-ux610fux5473ux7740ux4ec0ux4e48}

TCP 可以保证 byte 顺序，却不知道这段 byte 是域名查询、HTTP
request、邮件命令还是文件内容。Application protocol 必须定义：

\[
\boxed{\begin{gathered}
\text{Application object}\to\text{Name/identifier}\to\text{Operation}\\
\to\text{Message representation}\to\text{Response/state transition}
\end{gathered}}
\]

它不仅规定报文格式，还规定动作的语义与交互顺序。相同 transport
可以承载完全不同服务；相同服务语义也可以演进出不同 transport bindings。

\section{应用架构：谁长期持有状态}\label{2-ux5e94ux7528ux67b6ux6784ux8c01ux957fux671fux6301ux6709ux72b6ux6001}

\subsection{Client/Server}\label{21-clientserver}

server 通常在稳定地址/名字上等待请求并持有共享服务状态；client
主动发起交互。集中服务便于管理、一致性和发现，代价是容量、故障和运营中心化。

\subsection{Peer-to-Peer}\label{22-peer-to-peer}

peer 同时消费和提供资源，容量可随参与者增加，但发现、可用性、信任、NAT
穿越和一致性更复杂。

C/S 与 P2P 是状态与职责的组织方式，不等同于
TCP/UDP，也不由``是否有服务器进程''一个表象决定。

\section{DNS：把一个集中表拆成层次化授权系统}\label{3-dnsux628aux4e00ux4e2aux96c6ux4e2dux8868ux62c6ux6210ux5c42ux6b21ux5316ux6388ux6743ux7cfbux7edf}

\subsection{\texorpdfstring{为什么单一 \texttt{hosts}
文件失败}{为什么单一 hosts 文件失败}}\label{31-ux4e3aux4ec0ux4e48ux5355ux4e00-hosts-ux6587ux4ef6ux5931ux8d25}

全网共享一份名字表会遇到：更新冲突、分发延迟、单点容量与管理边界。DNS
的生成链是：

\[
\boxed{\begin{gathered}
\text{Global flat file fails}\to\text{Hierarchical namespace}\to\text{Delegated zones}\\
\to\text{Distributed authoritative servers}\to\text{Resolver + cache}
\end{gathered}}
\]

Domain namespace 是名字树；zone 是某个 authority
实际管理的数据边界。树的结构与管理分区相关，但 domain
不等于一台主机或一个物理网络。

\subsection{四类对象}\label{32-ux56dbux7c7bux5bf9ux8c61}

\begin{itemize}
\tightlist
\item
  stub resolver：应用/OS 发起查询的本地接口；
\item
  recursive resolver：代用户完成后续查询并缓存；
\item
  authoritative name server：对某 zone 的数据负责；
\item
  root/TLD/other delegated servers：通过 referral 把查询逐步引向
  authority。
\end{itemize}

\subsection{一次 cache miss
的解析轨迹}\label{33-ux4e00ux6b21-cache-miss-ux7684ux89e3ux6790ux8f68ux8ff9}

\begin{lstlisting}
application asks local resolver
-> recursive resolver checks cache
-> root referral identifies TLD servers
-> TLD referral identifies authoritative servers
-> authoritative answer returns records
-> resolver caches answer/referrals under TTL
-> result returns to application
\end{lstlisting}

迭代查询返回``我不知道最终答案，但下一步问谁''；递归查询把继续查找的责任交给被询问方。实际
client-to-recursive 常请求递归，resolver 对 authority 链执行迭代式追踪。

\subsection{Cache
的正确性边界}\label{34-cache-ux7684ux6b63ux786eux6027ux8fb9ux754c}

cache 用旧信息换低延迟和低查询负载。TTL 限制可复用时间，但 DNS
不是瞬时全局一致系统：更新在 TTL 生命周期内可能与缓存旧值并存。

DNS 并非只能用 UDP。经典普通查询常用 UDP；响应截断、zone transfer
等情形可使用 TCP，现代部署还存在加密传输。408
做题要服从题设，不能把``DNS=UDP''写成协议语义。

\section{HTTP：操作资源的统一语义}\label{4-httpux64cdux4f5cux8d44ux6e90ux7684ux7edfux4e00ux8bedux4e49}

\subsection{WWW：分布式超文档系统}

WWW 不是某一个 protocol 或“所有 Internet 服务”的同义词。它至少包含：browser/client、Web server、由 URL/URI 标识的 resource、HTML 等 document/representation，以及把文档与其他资源连接起来的 hyperlink。HTTP 定义 client 与 server 操作这些资源的请求/响应语义；DNS 可把 URL 中的 host name 解析为 IP；transport 再承载 HTTP message。

\[
\boxed{\text{URL identifies resource}\to\text{DNS locates host}\to
\text{HTTP requests representation}\to\text{HTML links reveal dependencies}}
\]

“网页”常是一个主文档与图片、样式、脚本等多个对象组成的依赖图，不是一条 HTTP response 必然完整返回的单文件。URL 提供标识与定位语法，HTML 提供超文本表示，HTTP 提供操作/传输语义；三者不能互相替代。

HTTP 的核心对象不是``网页文件''，而是由 URI 标识的
resource。消息传输的是 resource 的 representation，method 表达 client
希望对目标执行的语义。

\[
\boxed{
\text{Target resource}
\to \text{Method semantics}
\to \text{Request metadata/content}
\to \text{Status}
\to \text{Representation/metadata}
}
\]

\subsection{Resource 不等于
representation}\label{41-resource-ux4e0dux7b49ux4e8e-representation}

同一 resource 可根据时间、语言、编码或内容协商产生不同
representation。response body 是某次选定表示，不等于资源本身的全部身份。

\subsection{Method
属性改变重试决策}\label{42-method-ux5c5eux6027ux6539ux53d8ux91cdux8bd5ux51b3ux7b56}

\begin{itemize}
\tightlist
\item
  safe：语义上只读，例如 GET/HEAD；
\item
  idempotent：重复执行与执行一次具有相同预期 effect，例如 PUT、DELETE
  以及 safe methods；
\item
  POST 通常不假定 idempotent。
\end{itemize}

这些属性不是礼貌标签，而是 cache、proxy、client
在失败后能否自动重试的重要决策依据。

\subsection{Stateless
不等于``服务没有状态''}\label{43-stateless-ux4e0dux7b49ux4e8eux670dux52a1ux6ca1ux6709ux72b6ux6001}

HTTP stateless 表示每个 request 的语义不依赖 server 必须记住此前
protocol request context。应用仍可通过 cookie、token、database 和
server-side session 建立业务状态。

\subsection{语义与连接版本分离}\label{44-ux8bedux4e49ux4e0eux8fdeux63a5ux7248ux672cux5206ux79bb}

\begin{itemize}
\tightlist
\item
  HTTP/1.1 和 HTTP/2 通常在 TCP 上承载；
\item
  HTTP/3 使用 QUIC over UDP；
\item
  核心 method、status、representation 和 cache semantics
  在版本间保持共同定义。
\end{itemize}

因此``HTTP 使用 TCP''是特定版本/部署的 transport mapping，不是 HTTP
resource semantics 本身。

\subsection{Transport endpoint 不等于 HTTP authority：同一 IP:port 可以承载多个名字空间}

TCP 的四元组只把字节流交给某个 server endpoint；HTTP 仍需要知道当前请求属于哪个应用名字。对 HTTP/1.1，\code{Host} 请求首部携带目标 host/authority；HTTP/2/3 使用等价的 \code{:authority} 语义。因此两个不同域名即使 DNS 最终解析到同一 server IP、并使用相同 TCP port，也可以由同一服务进程按 application authority 选择不同 virtual host / resource namespace。

\[
\boxed{\text{Domain/Authority}\neq\text{IP address}\neq\text{TCP endpoint}}
\]

这里不能反向推断：共享 IP/port 不表示两个域名是同一应用资源；不同域名也不要求必须分配不同 IP 或不同 port。题目若问“同一服务器怎样区分两个基于名称的站点”，应在 HTTP message semantics 中找 Host/authority，而不是强行修改传输层标识。

\section{Web
性能从依赖图生成}\label{5-web-ux6027ux80fdux4eceux4f9dux8d56ux56feux751fux6210}

取得一个 URL 可能依次触发：DNS、transport/security establishment、HTTP
request/response、HTML parsing、更多 subresource requests。

总时间不能只数``几个 RTT''，必须先画依赖关系：哪些步骤串行、哪些
connection 可复用、哪些 object 可并行、哪些命中 cache。

Persistent connection 减少重复建立连接；multiplexing 允许多个
request/response 共享连接并发推进；cache 用 freshness/validation
状态减少传输。三者解决的瓶颈不同。

\section{FTP：控制状态与数据传输分离}\label{6-ftpux63a7ux5236ux72b6ux6001ux4e0eux6570ux636eux4f20ux8f93ux5206ux79bb}

FTP 的母问题是跨异构系统可靠操作和传输文件。它使用持久 control
connection 交换命令与会话状态，并为目录/文件数据建立独立 data
connection。

这种 out-of-band control
让命令不被大文件数据阻塞，也带来额外连接管理、主动/被动模式和
NAT/firewall 适配成本。FTP 的两条 TCP connections 是应用层设计，不是 TCP
自动生成的``控制通道''。

\section{电子邮件：提交、转发、存储、读取不是一个协议}\label{7-ux7535ux5b50ux90aeux4ef6ux63d0ux4ea4ux8f6cux53d1ux5b58ux50a8ux8bfbux53d6ux4e0dux662fux4e00ux4e2aux534fux8bae}

邮件系统至少有：user agent、mail server、message transfer 和 mailbox
access。

\begin{lstlisting}
sender user agent
-> submission / SMTP transfer
-> one or more mail servers
-> recipient mailbox
-> POP3 or IMAP access
-> recipient user agent
\end{lstlisting}

\begin{itemize}
\tightlist
\item
  SMTP 负责提交/服务器间 push 式传送；
\item
  POP3 提供较简单的 mailbox download/access 模型；
\item
  IMAP 更强调服务器端 mailbox 状态与多设备同步；
\item
  MIME 扩展 message representation，使非 ASCII
  内容和附件可被邮件体系承载。
\end{itemize}

``发邮件''和``收邮件''必须分开建模，不能把所有过程归给 SMTP。

\section{应用模型核验层：C/S 与 P2P 的可计算边界}

C/S 与 P2P 首先是角色、状态与资源位置的模型，不是 TCP/UDP 的同义词。若服务器以 $u_s$ 上传一个大小为 $F$ 的文件给 $N$ 个客户，客户最低下载速率为 $d_{min}$，在忽略核心拥塞的带宽下界模型中：
\[
  D_{C/S}=\max\left\{\frac{NF}{u_s},\frac{F}{d_{min}}\right\}.
\]
服务器上行要产出 $N$ 份完整副本，所以第一项随 $N$ 线性增长。

若对等方 $i$ 也能以上传速率 $u_i$ 分享文件，则
\[
  D_{P2P}=\max\left\{\frac{F}{u_s},\frac{F}{d_{min}},
  \frac{NF}{u_s+\sum_i u_i}\right\}.
\]
第三项的分子、分母同时随参与者增加；这叫 self-scalability，但不等于“人越多单个下载越快”。若所有 peer 都只下载、$u_i=0$，P2P 的供给项退化为 C/S 的服务器瓶颈。文件分块让新 peer 尚未下载完整文件时就能开始上传，是自扩展性真正落地的条件。

资源定位还要独立建模：集中目录（如集中 tracker）可以只集中“谁拥有”，文件仍由 peer 直接传；全分布式洪泛降低单点但增加查询开销；分块/邻居交换让定位与传输都分散。看到“P2P”不能直接推断没有任何中心辅助节点。

\section{DNS 核验层：名字树、授权区与记录类型}

一个 label 最多 63 个字符，完整域名文本总长通常不超过 255 个字符，大小写不区分；这些是名字语法边界，不是 IP 地址或子网掩码。domain 是名字树中的节点/子树，zone 是某个 authoritative server 实际拥有的管理数据边界，二者不能互换。

\begin{handbooktable}{L{2.4cm} L{2.5cm} Y}
\toprule
记录 & 指向/用途 & 题面判据\\
\midrule
A & 主机名到 IPv4 地址 & 返回 32-bit address\\
AAAA & 主机名到 IPv6 地址 & 返回 128-bit address\\
NS & zone 的权威服务器 & referral/委派链的下一责任者\\
CNAME & 别名到 canonical name & 还需继续解析目标名\\
MX & 域的邮件交换服务器及优先级 & SMTP relay 选择下一跳\\
PTR & 反向命名（地址到名称） & reverse lookup\\
\bottomrule
\end{handbooktable}

root server、TLD server、authoritative server 与 recursive resolver 的 Owner 不同。root/TLD 通常返回 referral（NS 与必要 glue），authoritative server 才对自己 zone 内最终 RR 负责；recursive resolver 代表 stub 承担继续查询、缓存和错误返回。迭代查询的响应把“下一步问谁”交还给 caller，递归查询则把后续责任交给被询问方。

一次 cache miss 的可审计顺序是：stub 发送 name/type/class/transaction ID；resolver 查本地 cache；若未命中，依次询问 root、TLD、authoritative；校验 response identity 与 RR；按 positive 或 negative TTL 缓存；向 stub 返回 answer 或错误。缓存提高性能却产生旧值窗口，TTL 到期是可复用期限，不是全网同时更新的保证。NXDOMAIN/negative answer 也可按其缓存语义保留，不能把“没查到”写成永远不存在。

普通 DNS query 常走 UDP 53；响应截断、zone transfer 或题设要求可靠字节流时可使用 TCP 53。现代加密 DNS 是 transport deployment 变化，不改变 hierarchy、authority 与 RR 语义。看到 `DNS=UDP` 只能说明常见承载，不能替代上述流程。

\section{HTTP 核验层：报文语法、方法属性与连接代价}

HTTP/1.x 请求/响应的语义骨架可以写成：
\[
\begin{aligned}
\text{request line:}&\quad \text{method + target + version},\\
&\quad \text{headers} + \text{blank line} + [\text{optional body}],\\
\text{response:}&\quad \text{version + status code/reason},\\
&\quad \text{headers} + \text{blank line} + [\text{representation}].
\end{aligned}
\]
常见方法的核心不是英文动词而是重试与缓存属性：GET/HEAD 通常 safe；PUT/DELETE 语义上 idempotent；POST 通常用于创建/触发处理且不能默认安全重试。状态码可按类别审计：1xx informational、2xx success、3xx redirection/cache validation、4xx client request、5xx server failure；状态码不是 TCP ACK，也不代表业务最终完成。

HTTP stateless 只说明每个 request 的协议语义不要求 server 保存上一次 request context；cookie、token、session/database 仍可维护业务状态。resource 是被 URI 标识的对象，representation 是某次按语言、编码或时间协商出的可传输表示；缓存控制必须针对 representation freshness/validation，而不能把 body 当永久资源身份。

性能题先分连接代价与消息代价。非持久 HTTP/1.0 的简单模型中，DNS 命中且无安全握手时，建立 TCP 约 1 RTT、发送 request 并得到首字节再约 1 RTT，之后再加对象传输时间；持久连接复用后可省后续建连 RTT；HTTP/1.1 pipeline、HTTP/2 multiplexing 与 HTTP/3/QUIC 改变并发与传输映射，但不改变 resource/method/status 语义。具体 RTT 公式必须按题设是否串行、缓存是否命中、对象是否并行来画依赖 DAG，不能套一个固定“网页 = 2 RTT”。

代理/Web cache 可能在 client 与 origin server 之间终止一段连接并用自己的 cache 回答；conditional GET 通过 validator（如 ETag/Last-Modified）让 origin 判断 304 Not Modified，减少 representation 重传。缓存命中、连接复用、multiplexing 是三个不同的优化点。

\section{邮件与 FTP 核验层：状态轨迹必须拆开}

邮件是 store-and-forward，而不是 sender 到 recipient user agent 的单条实时连接：
\[
\text{UA submission}\to\text{SMTP relay hop(s)}\to\text{recipient mailbox}\to\text{POP3/IMAP access}.
\]
SMTP 常用 TCP 25 做服务器间 relay（提交服务的端口与题设可能另给），采用 push；典型命令轨迹包含 220 greeting、HELO/EHLO、MAIL FROM、一个或多个 RCPT TO、DATA、单独一行句点结束、QUIT/221。SMTP reply 表示下一 hop 接受/拒绝，不表示收件人已经阅读。

邮件 content 与 envelope 分离：`To:`、`Cc:`、`Subject:` 属于内容 header，`MAIL FROM/RCPT TO` 是 SMTP envelope 命令；Bcc 可只存在 envelope 而不出现在可见 header。MIME 是内容扩充，不是独立传输协议：它用 Content-Type/Transfer-Encoding 等字段把非 ASCII 或附件编码成 SMTP 可承载的形式；ESMTP 是对 SMTP 命令/能力的扩展，二者不能混写。

POP3 是较简单的 pull/download 模型，常用 TCP 110；可下载并删除或保留副本，但本地文件夹/已读状态同步有限。IMAP 常用 TCP 143，更强调 mailbox 保留在 server、文件夹与多设备状态同步、按需下载部分内容。二者都不是 SMTP 的“接收阶段”同义词。

FTP 的 control/data 分离是应用层设计：control TCP 21 在会话期持续，传命令与认证；每次 LIST/RETR 等目录或文件数据操作另建 data TCP。主动模式由 server 从约定端口 20 连接 client 提供的 endpoint，被动模式由 server 通过 control 告知随机 data port、client 主动连接，后者更适合 NAT/firewall。一次会话若有 $m$ 次目录/文件数据传输，典型连接总数为 $1+m$，但同时存在通常只有 control 加当前一条 data。数据完成后 data state 释放，control state 继续保留。

FTP 传“副本”而非像 NFS 那样透明联机访问；ASCII 模式可做换行转换，binary/image 模式逐字节传输。TFTP 的 UDP/停止等待与小型启动用途属于扩展对照，不提升为本册 FTP canonical mechanism。

\section{408 应用协议流程：消息、状态与停止条件}

\subsection{DNS：递归责任与迭代 referral 分开}

\begin{figure}[htbp]
  \centering
  \resizebox{0.96\linewidth}{!}{\providecolor{themebg}{HTML}{FAFAF7}
\providecolor{themefg}{HTML}{111111}
\providecolor{themegray}{HTML}{666666}
\providecolor{themecurve}{HTML}{1D63B8}
\providecolor{themealert}{HTML}{C53030}
\providecolor{themeamber}{HTML}{B86E00}
\providecolor{themegreen}{HTML}{15803D}
\begin{tikzpicture}[font=\scriptsize,color=themefg,text=themefg,>=Stealth]
\tikzset{
  panel/.style={draw=themegray!50,rounded corners=6pt,fill=themefg!2!themebg,inner sep=8pt},
  box_cli/.style={draw=themecurve,rounded corners=3pt,minimum width=2.4cm,minimum height=.75cm,align=center,fill=themecurve!10!themebg,font=\scriptsize\bfseries,text=themecurve},
  box_dns/.style={draw=themeamber,rounded corners=3pt,minimum width=2.6cm,minimum height=.75cm,align=center,fill=themeamber!10!themebg,font=\scriptsize\bfseries,text=themeamber},
  box_srv/.style={draw=themegreen,rounded corners=3pt,minimum width=2.6cm,minimum height=.75cm,align=center,fill=themegreen!10!themebg,font=\scriptsize\bfseries,text=themegreen},
  lbl/.style={font=\tiny\bfseries,fill=themebg,inner sep=1.5pt,rounded corners=1pt},
  flow_q/.style={->,line width=1.1pt,themecurve},
  flow_r/.style={->,line width=1.1pt,themeamber,dashed},
  note/.style={rounded corners=4pt,draw=themegray!70,fill=themefg!3!themebg,inner xsep=8pt,inner ysep=7pt,align=left}
}

% ── 标题与例题说明 ──
\node[anchor=west,font=\bfseries\Large] at (0,8.8) {DNS 域名系统：递归查询（客户端到本地域名服务器）与迭代查询（分级权威）};
\node[anchor=west,text=themegray,text width=17.5cm,align=left] at (0,8.2)
  {解析场景：主机查询 \texttt{www.abc.com}。主机到本地域名服务器采用递归查询；本地域名服务器向上级采用迭代查询。};

% ── 拓扑节点定义 ──
\node[box_cli] (host) at (1.4,5.0) {客户端主机\\(Client)};
\node[box_dns] (local) at (5.2,5.0) {本地域名服务器\\(Local DNS)};

\node[box_srv] (root) at (14.0,7.3) {根域名服务器\\(Root DNS)};
\node[box_srv] (tld) at (14.0,5.0) {顶级域名服务器\\(.com TLD)};
\node[box_srv] (auth) at (14.0,2.7) {权威域名服务器\\(abc.com Auth)};

% ── 递归查询 (主机 <-> Local DNS) ──
\draw[flow_q] (host.north east) to[out=25,in=155] node[above=1pt,lbl,text=themecurve] {1. 递归查询: \texttt{www.abc.com}} (local.north west);
\draw[flow_r] (local.south west) to[out=-155,in=-25] node[below=1pt,lbl,text=themegreen] {8. 最终返回 IP 结果} (host.south east);

% ── 迭代查询 (Local DNS <-> 上级服务器) ──
% 步骤 2 & 3: 根 (y 范围 6.0 .. 7.5)
\draw[flow_q] (6.5,5.6) to[out=30,in=180] node[pos=0.45,above=2pt,lbl,text=themecurve] {2. 查询根} (root.west);
\draw[flow_r] (root.south west) to[out=-155,in=15] node[pos=0.6,below=2pt,lbl,text=themeamber] {3. 推荐 .com TLD} (6.5,5.3);

% 步骤 4 & 5: TLD (y 范围 4.7 .. 5.2)
\draw[flow_q] (6.5,5.0) -- node[pos=0.5,above=2pt,lbl,text=themecurve] {4. 查询 .com} (12.7,5.0);
\draw[flow_r] (12.7,4.6) -- node[pos=0.5,below=2pt,lbl,text=themeamber] {5. 推荐 abc.com 权威} (6.5,4.6);

% 步骤 6 & 7: 权威 (y 范围 2.5 .. 4.5)
\draw[flow_q] (6.5,4.3) to[out=-15,in=160] node[pos=0.55,above=2pt,lbl,text=themecurve] {6. 查询权威} (auth.north west);
\draw[flow_r] (auth.west) to[out=-175,in=-30] node[pos=0.45,below=2pt,lbl,text=themegreen] {7. 返回目标 IP 地址} (6.5,4.0);

% ── 本地缓存机制 ──
\node[draw=themegreen,rounded corners=3pt,fill=themegreen!8!themebg,font=\tiny\bfseries,text=themegreen,anchor=north,align=center] at (5.2,3.4)
  {DNS 高速缓存 (TTL 倒计时)\\命中直接返回 8，无需发起 2..7};

% ── 底部总结与误读警示 ──
\node[note,anchor=west,text width=17.5cm] at (0,1.0)
  {\textbf{DNS 递归 vs 迭代查询核心区别：}\\
   1. \textbf{递归查询 (Recursive)}：被询问服务器承担“包办到底”责任，自身去查询其他服务器并最终返回 IP（客户端 $\leftrightarrow$ 本地 DNS）；\\
   2. \textbf{迭代查询 (Iterative)}：被询问服务器仅返回“下一级服务器地址”，由请求方自己再去向下一级发起查询（本地 DNS $\leftrightarrow$ 根/TLD/权威）；\\
   3. \textbf{传输层承载}：常规域名解析使用 \textbf{UDP 53}（开销小、响应快）；区域传送 (Zone Transfer) 使用 \textbf{TCP 53}（保证主从同步可靠性）。};

\node[anchor=west,font=\scriptsize\bfseries,text=themealert,text width=17.5cm,align=left] at (0,-0.15)
  {禁止误读：根域名服务器通常不直接返回目标主机的 IP，而是返回负责顶级域名 (.com, .cn) 的 TLD 服务器地址。};
\end{tikzpicture}
}
  \caption{DNS 域名系统：递归查询与迭代查询分级权威解析}
  \label{fig:dns-query-modes}
\end{figure}


stub 向 recursive resolver 发送 query（name、type、class、transaction ID 等）；resolver 先查 cache。未命中时，它依次向 root、TLD 和 authoritative server 发迭代查询，referral 通过 NS 及必要的 glue/address 信息指出下一位 authority；最终 answer 或 negative result 按 TTL 缓存并返回 stub。收到 response 时必须核对 query identity/匹配字段；TTL 到期后旧 cache 不再直接回答。UDP 是经典普通查询的常见承载，截断或 zone transfer 等分支可转 TCP，不能把承载写成唯一语义。

\subsection{HTTP：请求语义与连接时延分开}

一次交互至少推进：
\[
\boxed{\text{method + target + fields + optional content}
\to\text{server processing}
\to\text{status + fields + optional representation}}.
\]
性能题先画依赖：DNS 是否命中、transport 是否已有连接、连接能否复用、对象是否串行/并行、内容大小和链路速率。Persistent connection 减少重复建连，cache 减少请求/传输，multiplexing 改变并发关系；三者不能用同一个“少一个 RTT”口诀替代。

\subsection{FTP：控制连接决定数据连接怎样建立}

client 先建立持久 control connection 并完成认证/命令。需要目录或文件数据时，主动模式由 server 从约定 data port 连接 client 提供的 endpoint；被动模式由 server 提供监听 endpoint，client 主动连接。一次 data transfer 完成后数据连接可关闭，control state 继续存在。端口常数和命令细节以题设为准，稳定边界是 control/data 两条状态轨迹。

\subsection{Mail：提交、relay 与 access 三段}

sender user agent 把 message 提交给 SMTP server；服务器通过 DNS MX 等信息选择下一跳并用 SMTP relay 推送到 recipient mail server；message 进入 mailbox 后，recipient 再用 POP3/IMAP 等访问协议读取。SMTP reply 驱动提交/转发状态，MIME 只定义可承载的内容表示；它们都不拥有 mailbox access。停止条件应分别写“已被下一 SMTP hop 接受”和“用户已从 mailbox 访问”，不能合并成一次端到端 ACK。

\begin{warnbox}{Owner 边界}
DHCP 虽常列在应用层协议表中，但本项目由 NET04 Own其配置生命周期；NET-I01 只调用 DNS、TCP 和 HTTP 的输出，不重复本节报文机制。
\end{warnbox}

\section{怎样选择 transport
service}\label{8-ux600eux6837ux9009ux62e9-transport-service}

应用选择的不是协议名偏好，而是约束组合：

\begin{longtable}[]{@{}ll@{}}
\toprule\noalign{}
应用需求 & 需要观察的 transport 能力 \\
\midrule\noalign{}
\endhead
\bottomrule\noalign{}
\endlastfoot
不能丢 byte、允许额外延迟 & reliable ordered stream \\
保留 message boundary、自定义恢复 & datagram \\
极低启动延迟 & connection setup cost \\
多路流并发且减少队头阻塞 & multiplexing semantics \\
广播 bootstrap & local datagram/broadcast support \\
安全身份与机密性 & security layer/secure transport \\
\end{longtable}

``实时应用一定用 UDP''\,``可靠应用一定用 TCP''都过度简化。应用可在 UDP
之上实现可靠/拥塞控制，或为穿越部署环境选择
TCP；判断必须回到服务需求与实现责任。

\section{概念边界}\label{9-ux6982ux5ff5ux8fb9ux754c}

\begin{handbooktable}{L{2.1cm} C{0.35cm} L{2.25cm} Y Y}
\toprule\noalign{}
概念 A & ≠ & 概念 B & 真正区别与题目信号 & 混淆后果 \\
\midrule\noalign{}
Domain name & ≠ & IP address & 人/组织命名与网络定位；DNS 维护映射 & 把
DNS 当逐跳路由 \\
Domain & ≠ & Zone & 名字树节点/子树与 authority 管理边界 &
授权与层次结构混乱 \\
Recursive query & ≠ & Iterative query & 被询问方完成解析；被询问方返回
referral & DNS 报文轨迹画错 \\
Resource & ≠ & Representation & 被标识对象与某次可传输表示 &
缓存和内容协商理解错误 \\
HTTP stateless & ≠ & Application has no state &
协议请求独立与业务状态存在可并存 & 认为登录/cookie 违反 HTTP \\
HTTP semantics & ≠ & HTTP transport version & method/status/resource
与具体连接承载分离 & 把 HTTP/3 误判为非 HTTP \\
SMTP & ≠ & POP3/IMAP & 邮件传送与 mailbox access &
邮件全流程协议归属错误 \\
FTP control connection & ≠ & FTP data connection &
命令状态与内容传输分离 & 端口/连接数量推演错误 \\
\bottomrule
\end{handbooktable}

\section{做题调用协议}\label{10-ux505aux9898ux8c03ux7528ux534fux8bae}

\begin{enumerate}
\def\labelenumi{\arabic{enumi}.}
\tightlist
\item
  写应用对象：name、resource、file、message 还是 lease/config；
\item
  写谁持有 authority/state，谁只是 cache/client；
\item
  画 request/response 或 push/pull 的事件方向；
\item
  DNS 题区分 recursion、iteration、authority、cache TTL；
\item
  HTTP 题区分 resource、representation、method 和 connection；
\item
  邮件题拆成 submission/transfer/storage/access；
\item
  性能题画依赖 DAG，再数不可并行的 RTT/transfer；
\item
  最后检查 transport 选择是否来自应用约束而非口诀。
\end{enumerate}

\section{贯穿母例：输入 URL 后第一个 HTTP request
为什么还不能立刻发}\label{11-ux8d2fux7a7fux6bcdux4f8bux8f93ux5165-url-ux540eux7b2cux4e00ux4e2a-http-request-ux4e3aux4ec0ux4e48ux8fd8ux4e0dux80fdux7acbux523bux53d1}

\begin{lstlisting}
URL identifies scheme + authority + target
-> resolver needs server IP (unless cached)
-> IP forwarding needs next hop and local MAC (unless cached)
-> transport/security context may need establishment
-> HTTP request names resource and method
-> response carries status + selected representation
-> representation may reveal more resource dependencies
\end{lstlisting}

URL、domain、IP、port、MAC
分别服务不同作用域。把它们压成一句``浏览器访问服务器''，会隐藏几乎所有可诊断的状态交接。

\section{高频 First
Divergence}\label{12-ux9ad8ux9891-first-divergence}

\begin{itemize}
\tightlist
\item
  把 DNS 说成从 root 一次返回最终 IP：漏掉 delegation/referral；
\item
  认为 cache 命中永远是最新真相：忽略 TTL 与弱一致时间窗；
\item
  写``DNS 用 UDP，HTTP 用 TCP''作为完整机制：把常见承载当唯一语义；
\item
  把 HTTP stateless 理解成网站不能登录：混淆 protocol context 与
  application state；
\item
  把 response body 等同 resource：混淆对象与 representation；
\item
  用 SMTP 从 mailbox 下载邮件：没有分开 transfer 与 access。
\end{itemize}

\section{一页压缩与复原问题}\label{13-ux4e00ux9875ux538bux7f29ux4e0eux590dux539fux95eeux9898}

\[
\boxed{\begin{gathered}
\text{Name an application object}\to\text{Choose operation semantics}\to
\text{Represent in messages}\\
\to\text{Map onto transport}\to\text{Evolve application state}
\end{gathered}}
\]

\begin{enumerate}
\def\labelenumi{\arabic{enumi}.}
\tightlist
\item
  DNS 怎样用 delegation 替代全球单表？
\item
  cache 为什么同时提高性能并引入旧信息窗口？
\item
  HTTP resource 与 representation 为什么必须分开？
\item
  stateless protocol 怎样支撑有状态登录业务？
\item
  SMTP、POP3/IMAP 和 MIME 分别拥有邮件过程的哪一段？
\end{enumerate}

\section{来源与校正说明}\label{14-ux6765ux6e90ux4e0eux6821ux6b63ux8bf4ux660e}

\begin{itemize}
\tightlist
\item
  归档笔记《应用层-FTP,SMTP,POP3,DHCP》提供传统应用协议覆盖；《应用层-DNS》《应用层-WWW与HTTP》为空文件，未被当作已有正文；
\item
  DNS 的 hierarchy、delegation、resolver、authority 与 cache 边界依据
  \href{https://www.rfc-editor.org/rfc/rfc1034.html}{RFC 1034} 校正；
\item
  HTTP 的 stateless semantics、resource/representation
  和跨版本语义边界依据
  \href{https://www.rfc-editor.org/rfc/rfc9110.html}{RFC 9110} 校正；
\item
  DHCP 已按当前 Ownership 移入 IP Topic，不在应用册重复拥有。
\end{itemize}
\end{document}
